\subsection{Problem 11. Riemann sums and a parameter}
Let $f(x)=x^{2}+2x$ on the interval $[0, a]$, where $a>0$ is a constant.

\begin{enumerate}[label=(\alph*)]
    \item Construct the right-endpoint Riemann sum $R_{n}$ for $\displaystyle\int_{0}^{a}f(x)dx$ using $n$ equal subintervals. Express $R_{n}$ explicitly as a sum.
    \item Rewrite $R_{n}$ using sigma notation and known formulas for $\sum_{k=1}^{n}k$ and $\sum_{k=1}^{n}k^{2}$.
    \item Compute $\lim_{n\rightarrow\infty}R_{n}$ and verify that the result matches the antiderivative computation of $\displaystyle\int_{0}^{a}f(x)dx$.
    \item Interpret your final expression as an area and discuss how it depends on the parameter $a$. What happens if you double $a$?
\end{enumerate}

\subsection*{Strategy}
\begin{enumerate}[label = (\alph*)]
    \item To construct the right-endpoint Riemann sum $R_n$, we first construct the $n$ equal subintervals. We then calculate the area for each rectangle in the corresponding subinterval and express $R_n$ as the sum of these areas.
    \item To rewrite $R_n$, we use mathematical operations to find out the formula of $R_n$.
    \item We first compute $\lim_{n\to\infty} R_{n}$. Next, we calculate the antiderivative computation of $\int_0^af(x)dx$ and illustrate that the values computed are the same.
    \item We express the final expression geometrically and show the relationship between the area and the parameter a.
\end{enumerate}
\subsection*{Solution}

\textbf{(a) Construction and expression of $\mathbf{R_n}$}
    \begin{itemize}
    \item \textbf{Construction:} \\
    To start with, we construct the $n$ equal subintervals. Since there are $n$ equal subintervals in the interval $[0, a]$, the width of each rectangle in the corresponding subinterval is: 
    \begin{equation*}
        \Delta x = \frac{a - 0}{n} = \frac{a}{n}
    \end{equation*}
    We denote $x_i$ as the endpoint of the subintervals. The $n$ equal subintervals can be displayed as: 
    \[[x_0, x_1], [x_1, x_2],...,[x_{n-1}, x_n]\]
    \centerline {Where: \(x_i = x_0 + i\Delta x = \frac{ia}{n}\)} \\
    Since we are dealing with right-endpoint, the height of the rectangle in the subinterval $[x_{i-1}, x_i]$ is $f(x_i)$. \\
    Therefore, we achieve the area of the $i^{th}$ rectangle as: \(f(x_i)\Delta x\). \\
    \item \textbf{Expression:} \\
    Taking the sum of all of the rectangle areas gives us: 
    \begin{equation*}
    \begin{split}
        R_n &= \sum_{i=1}^n f(x_i)\Delta x \\
        &= \sum_{i=1}^n (x_i^2 + 2x_i)\frac{a}{n} \\
        &= \sum_{i=1}^n [(\frac{ia}{n})^2 + 2\frac{ia}{n}]\frac{a}{n} \\
        &= \sum_{i=1}^n (\frac{i^{2}a^{3} + 2ia^{2}n}{n^3})
    \end{split}
    \end{equation*}
    Hence,
    \begin{center} \fbox{$\displaystyle R_n = \sum_{i=1}^n (\frac{i^{2}a^{3} + 2ia^{2}n}{n^3})$} \end{center}
    \end{itemize}

\textbf{(b) Rewriting $\mathbf{R_n}$} \\
    To rewrite $R_n$, we first notice that the formula of $R_n$ contains the constant \(\displaystyle\frac{a^2}{n^2}\), which can be put outside the sigma notation.\\
    We rewrite $R_n$ as: \[R_n = \frac{a^2}{n^2}\sum_{i=1}^n (\frac{i^{2}a}{n} + 2i)\]
    Let \[S = \sum_{i=1}^n (\frac{i^{2}a}{n} + 2i)\] 
    \begin{equation}
    \Longrightarrow R_n = \frac{a^2}{n^2}\cdot S
    \end{equation}
    We calculate $S$ :
    \begin{align*}
        S &= \sum_{i=1}^n (\frac{i^{2}a}{n} + 2i) \\
        &= \frac{a}{n}\sum_{i=1}^n i^2 + 2\sum_{i=1}^ni\\
    \end{align*}
    By applying the formulas:  \[\sum_{k=1}^n k = \frac{n(n+1)}{2}\quad; \quad\sum_{k=1}^n k^2 = \frac{n(n+1)(2n+1)}{6}\], we get: \\
    \begin{align*}
        S &= \frac{a}{n} \cdot \frac{n(n+1)(2n+1)}{6} + 2\cdot \frac{n(n+1)}{2} \\
        \iff S &= \frac{a(n+1)(2n+1)}{6} + n(n+1) 
    \end{align*}
    Substituting the formula of S into (3) gives us:
    
    \begin{center} \fbox{$\displaystyle R_n = \frac{a^3 (n+1)(2n+1)}{6n^2} + \frac{a^2 (n+1)}{n}$} \end{center}

\textbf{(c) Limit computation and verification:}
    \begin{itemize} 
    \item \textbf{Limit calculation:}
    \begin{align*}
        \lim_{n\to \infty} R_n &= \lim_{n\to \infty} [\frac{a^3 (n+1)(2n+1)}{6n^2} + \frac{a^2 (n+1)}{n} ] \\
        &= \lim_{n\to \infty} [\frac{a^3 (n+1)(2n+1) + a^2 6n(n+1)}{6n^2} ] \\
        &= \lim_{n\to \infty} [\frac{(2a^3 + 6a^2)n^2 + (3a^3 + 6a^2)n + a^3}{6n^2} ] \\
        &= \lim_{n\to \infty} [\frac{(2a^3 + 6a^2) + \frac{3a^3 + 6a^2}{n} + \frac{a^3}{n^2}}{6} ] \\
        \intertext{Using the limit property $\displaystyle\lim_{n \to \infty} \frac{c}{n^k} = 0$ for $k > 0$, we get: } \\
        \lim_{n\to \infty} R_n &= \lim_{n\to \infty} [\frac{(2a^3 + 6a^2) + 0 + 0}{6} ] \\
        &= \lim_{n\to \infty} (\frac{a^3}{3} + a^2) \\
        &= \frac{a^3}{3} + a^2 
    \end{align*}
    Thus, \begin{equation}
    \lim_{n\to \infty} R_n = \frac{a^3}{3} + a^2
    \end{equation}
    \item \textbf{Antiderivative computation:}
    \begin{align*}
        \int_0^af(x)dx &= \int_0^a (x^2 + 2x)dx \\
        &= \left. \frac{x^3}{3} + x^2 \right|_{x=0}^{x=a}\\
    \end{align*}
    \begin{equation}
        \Longleftrightarrow \int_0^af(x)dx  = \frac{a^3}{3} + a^2
    \end{equation}
    \item \textbf{Verification:}
    From (4) and (5), it can be seen that the Riemann sum limit matches the antiderivative result. Hence, the computation is verified.
    \end{itemize}
    
\textbf{(d) Interpretation and Parameter Dependence:}
    The final expression $A(a) = \frac{a^3}{3} + a^2$ represents the \textbf{exact area} under the curve $f(x)=x^2+2x$ from $x=0$ to $x=a$.
    
    \begin{itemize}
        \item \textbf{Dependence on $\mathbf{a}$:} The area $A(a)$ is a polynomial function of degree 3 in terms of the parameter $a$. This confirms that the area is highly sensitive to the width of the interval.
    
        \item \textbf{Doubling $\mathbf{a}$:} If we double the width of the interval from $a$ to $2a$, the new area is $A(2a)$:
        \[
        A(2a) = \frac{(2a)^3}{3} + (2a)^2 = \frac{8a^3}{3} + 4a^2
        \]
        Since $A(2a) \neq 2 \cdot A(a)$, doubling the width of the interval more than doubles the area due to the cubic and quadratic growth terms.
    \end{itemize}
\pagebreak